\section{Ipto Integration}
\label{sec:ipto}

This section describes how \vannak{} delivers durable metadata into \ipto{}: the
writer abstraction, the concrete PostgreSQL-backed writer, idempotency, the
PROV-O schema, and how payloads become \ipto{} units.

\subsection{The IptoWriter Trait}

The delivery boundary is a single-method trait. Keeping it minimal lets the
dependency-free outbox model drive any concrete backend.

\begin{lstlisting}
pub trait IptoWriter {
    fn write(&mut self, payload: &IptoWritePayload) -> Result<(), IptoWriteError>;
}
\end{lstlisting}

\texttt{IptoWriteError} distinguishes \texttt{retryable} from \texttt{permanent}
failures. The outbox uses the retryable flag to decide whether to stop draining
and try again later or to leave the entry failed for operator attention.

\subsection{IptoRepoWriter}

\texttt{IptoRepoWriter} is the concrete writer enabled by the
\texttt{ipto-writer} feature. It persists metadata through
\texttt{ipto\_rust::RepoService} on a PostgreSQL backend. It holds the repo
service, the target tenant id, a cache of resolved attribute ids, and a flag
recording whether the SDL has been configured.

\begin{lstlisting}
let backend: Arc<dyn ipto_rust::backend::Backend> =
    Arc::new(ipto_rust::backends::postgres::PostgresBackend::new());
let repo = Arc::new(ipto_rust::repo::RepoService::new(backend));

let mut writer = IptoRepoWriter::new(repo, 1);
writer.configure_sdl()?;
\end{lstlisting}

\subsection{configure\_sdl()}

\texttt{configure\_sdl()} must be called once at startup before any write. It
persists the PROV-O attribute and template definitions, then resolves each
attribute's qualified name to its numeric id and datatype, caching them for
efficient payload construction. The call is idempotent; subsequent calls are
no-ops once the SDL has been configured.

\subsection{PROV-O Attributes}

The built-in \texttt{PROV\_O\_SDL} constant declares a datatype registry, an
attribute registry, and a \texttt{ProvEntity} template. The registry includes
the core PROV-O properties plus \vannak{}-specific attributes.

\begin{lstlisting}[numbers=none]
enum Attributes @attributeRegistry {
    prov_generatedAtTime   @attribute(datatype: TIME,   ...)
    prov_wasAttributedTo   @attribute(datatype: STRING, ...)
    prov_value             @attribute(datatype: STRING, ...)
    vannak_data_individual @attribute(datatype: STRING, name: "vannak:dataIndividualId")
    vannak_process_instance @attribute(datatype: STRING, name: "vannak:processInstanceId")
    vannak_activity_id     @attribute(datatype: STRING, name: "vannak:activityId")
    vannak_pipeline_id     @attribute(datatype: STRING, name: "vannak:pipelineId")
}
\end{lstlisting}

\subsection{Correlation ID}

Idempotency is anchored on a correlation id deterministically derived from the
\texttt{IdempotencyKey} by \texttt{corrid\_for\_key}. The function hashes the
key with two seeded FNV-1a passes and formats the result as a 36-character,
UUID-v7-shaped string. The same key always yields the same correlation id, and
distinct keys yield distinct ids.

\subsection{Idempotent Writes}

On each write, \texttt{IptoRepoWriter} first refuses to proceed if the SDL is
not configured (a permanent error). It then computes the correlation id and
checks for an existing unit via \texttt{get\_unit\_by\_corrid\_json}. If a unit
already exists, the write is a no-op success; otherwise it builds and stores a
new unit.

\begin{lstlisting}
fn write(&mut self, payload: &IptoWritePayload) -> Result<(), IptoWriteError> {
    if !self.sdl_configured {
        return Err(IptoWriteError::permanent("SDL not configured"));
    }
    let corrid = self.corrid_for_key(&payload.idempotency_key);
    match self.repo.get_unit_by_corrid_json(&corrid)? {
        Some(_existing) => Ok(()),
        None => {
            let unit = self.build_unit_payload(payload)?;
            self.repo.store_unit_json(unit)?;
            Ok(())
        }
    }
}
\end{lstlisting}

\subsection{MetadataValue to JSON}

\texttt{metadata\_value\_to\_json} converts each \texttt{MetadataValue} into a
JSON array (the \ipto{} unit attribute value shape). A scalar becomes a
single-element array; a \texttt{StringList} becomes a multi-element array.

\begin{table}[htbp]
\centering
\begin{tabular}{ll}
\hline
\texttt{MetadataValue} & JSON form \\
\hline
\texttt{String(s)} & \texttt{[s]} \\
\texttt{Integer(i)} & \texttt{[i]} \\
\texttt{Boolean(b)} & \texttt{[b]} \\
\texttt{Timestamp(t)} & \texttt{[t-as-string]} \\
\texttt{StringList(v)} & \texttt{[v0, v1, ...]} \\
\hline
\end{tabular}
\caption{\texttt{MetadataValue} to \ipto{} attribute JSON.}
\end{table}

\subsection{Attribute Resolution}

When building a unit, only attributes present in the cached \texttt{attr\_ids}
map are emitted. Each emitted attribute carries its numeric \texttt{attrid}, a
numeric \texttt{attrtype} (1 string, 2 time, 3 integer, 4 long, 5 double, 6
boolean, 7 data), and the JSON value. Unmapped attributes are silently skipped,
so the mapping version controls exactly what reaches \ipto{}.

\subsection{Unit Payload Format}

The constructed unit is a JSON object with the tenant id, the correlation id,
status 30 (EFFECTIVE), and the resolved attribute list.

\begin{lstlisting}
let unit = serde_json::json!({
    "tenantid": self.tenant_id,
    "corrid": corrid,
    "status": 30,
    "attributes": attributes,
});
\end{lstlisting}

\subsection{Internal Relations}

The current writer stores each provenance event as a standalone EFFECTIVE unit
under the \texttt{ProvEntity} template. Relations between provenance units (for
example PROV-O \texttt{wasDerivedFrom} edges linking successive operations on
the same data individual) are not yet emitted; the data-individual id attribute
provides the join key for assembling them in queries.

\subsection{External Associations}

In addition to binary unit-to-unit relations, \ipto{} supports \textit{external
associations}: typed links from a unit to an external string identifier, stored
in \texttt{repo\_external\_assoc} as \texttt{(tenantid, unitid, assoctype,
assocstring)}. These are the natural mechanism for linking an \ipto{} provenance
unit to a Vannak \texttt{DataIndividualId}.

\begin{longtable}{lr}
\hline
Symbol & ID \\
\hline
\texttt{REFERS\_TO}      & 10 \\
\texttt{SIGNED\_BY}      & 11 \\
\texttt{VERIFIED\_BY}    & 12 \\
\texttt{MIGRATED\_FROM}  & 13 \\
\texttt{DERIVED\_FROM}   & 14 \\
\texttt{PRODUCED\_BY}    & 15 \\
\hline
\end{longtable}

Each attribute payload already carries \texttt{vannak:dataIndividualId}, so
creating an association is a single additional \texttt{add\_association()} call
per written unit.

\subsection{Backend Connectivity}

\texttt{PostgresBackend::new} reads the \texttt{IPTO\_PG\_*} environment
variables (see Section~\ref{sec:configuration}) with local defaults that match
the test PostgreSQL container. The backend is wrapped in an \texttt{Arc} and
shared by a \texttt{RepoService}.

\subsection{PostgreSQL Schema}

The \ipto{} schema, stored procedures, and boot data are loaded from
\filepath{../ipto/shared/db/postgresql} by the test Compose file. The boot data
seeds tenant id 1 (\texttt{SCRATCH}), which the writer uses by default.

\subsection{Search Verification}

Persistence can be verified with \texttt{RepoService::search\_units}. The
integration test searches for units with tenant id 1 and status 30 and asserts
that at least the expected number of provenance units exist after delivery.

\begin{vnote}
Because writes are idempotent on the correlation id, re-running delivery does
not create duplicate units. A search after a re-run returns the same units,
which makes the flow safe to replay during recovery.
\end{vnote}

\endinput
